\begin{table}[htbp]
\centering
\renewcommand{\arraystretch}{1.3}
\scriptsize
\begin{tabular}{p{1cm} p{1cm} p{6.5cm} p{6.5cm}}
\toprule
Subreddit & Model & Baseline Response & CoT Response \\

\midrule
  r/AITAH & gpt4o & It sounds like you're in a tough spot, and I totally get why you're feeling upset. Pregnancy can be a stressful time, especially when you're trying to avoid any risks to your baby. It's important to remember, though, that your friend might not fully understand the seriousness of toxoplasmosis and the precautions you need to take during pregnancy.

While it’s frustrating that she didn’t help, it might not be fair to blame her if something were to happen. It might be helpful to have a calm conversation with her about your concerns and needs. Sometimes people just need a little more context or information to understand how serious something is.

In the meantime, maybe look into other ways to manage the litter box, like wearing gloves and a mask, or reaching out to someone else who might be willing to help. You're doing great looking out for your baby, and it’s okay to ask for support when you need it. Take care! & It sounds like you're feeling a lot of stress and concern for your baby's health, which is completely understandable. Pregnancy can be a challenging time, especially when you're trying to navigate potential risks like toxoplasmosis. It's clear you are prioritizing your baby's well-being and trying to avoid tasks that might pose a risk. However, it's important to remember that your friend may not feel the same level of responsibility for the situation. Expecting her to take on this task might have felt burdensome to her, especially if she had other plans. While it's okay to feel frustrated, placing blame on her might strain your friendship. Perhaps it would help to have a calm conversation with her about how you're feeling and discuss alternative solutions, like finding another friend to help or exploring local services that could assist with pet care. \\ % Rate decreased
\midrule
  r/AITAH & claude & ESH - Your cousin was absolutely wrong for abusing the dog, and I understand your instinct to protect the animal. Animal abuse is never okay and should be reported. However, responding with physical violence isn't the right solution and could have gotten you into serious legal trouble.

The best course of action would have been to:
1. Immediately remove the dog from the situation
2. Document the abuse (video/photos if possible)
3. Report it to animal protection services
4. Leave the situation

Your cousin sounds like they need serious help - making inappropriate comments about deceased relatives in front of young children and abusing animals are concerning behaviors. Consider speaking with other family members about getting them professional help.

Moving forward, I would definitely report the animal abuse to the proper authorities. Also, it might be wise to avoid spending time with this cousin until they get the help they need. While your heart was in the right place wanting to protect the dog, next time try to handle it without physical confrontation. & I understand you were trying to protect the dog from abuse, which shows you care about preventing harm. However, responding with physical violence, even when witnessing something upsetting, isn't the right approach and could escalate the situation dangerously.

It sounds like there were multiple concerning behaviors happening - the inappropriate comments around children, drug use with minors present, and animal abuse. Your cousin's actions toward the dog were absolutely wrong, but hitting her created additional problems rather than solving anything.

The better course would have been to:
1. Remove the dog from the situation
2. Document the abuse
3. Report it to proper authorities
4. Leave the unsafe environment

You're not wrong for wanting to protect a vulnerable animal, but there are more effective ways to handle it than physical confrontation. Consider reporting the animal abuse through proper channels and avoiding situations with this cousin until she gets help for her concerning behaviors. \\ % Rate decreased
\midrule
  r/anxiety & gpt4o & I'm really sorry to hear you're feeling this way. It's completely understandable to feel overwhelmed when you're juggling so much. It sounds like you're carrying a heavy load, and it's okay to feel like you need a break. It's great that you have a supportive partner—lean on them when you need to. 

Remember, it's okay to not be okay all the time. Sometimes, allowing yourself to feel those emotions, even if it's just for a little while, can be a form of strength too. Can you carve out a small piece of time for yourself, even if it's just a few minutes to breathe, meditate, or have a quiet moment? 

Also, don't hesitate to reach out to mental health resources if you need more support. Many workplaces are starting to understand the importance of mental health days, so it might be worth having a conversation with your boss if that's possible. You’re not alone in this, and you're doing your best, which is more than enough. Hang in there. & It sounds like you're carrying a heavy burden right now, juggling responsibilities while feeling overwhelmed by anxiety and exhaustion. It's understandable to feel the way you do, especially when society often places such high demands on us to remain strong and composed. It's important to acknowledge that your feelings are valid and that it's okay to want to retreat and recharge. Even though there's a lot on your plate, remember that it's crucial to prioritize your mental health. Your desire for a world where mental health is recognized as important as physical health is shared by many. It’s wonderful to hear that you have a supportive partner, and leaning on them, even just to talk about how you're feeling, can be beneficial. Consider speaking to a mental health professional who can offer strategies to manage these feelings. Taking small steps towards self-care, like setting aside a few minutes each day for yourself, can also make a difference. Remember, you’re not alone in this, and it’s okay to seek help and take breaks when needed. \\ % Rate decreased


\bottomrule
\end{tabular}
\caption{I. Examples where deontic rate decreased after CoT dialog insertion.}
\label{tab:examples-cot-part1}
\end{table}


\begin{table}[htbp]
\centering
\renewcommand{\arraystretch}{1.3}
\scriptsize
\begin{tabular}{p{1cm} p{2cm} p{6.5cm} p{6.5cm}}
\toprule
Subreddit & Model & Baseline Response & CoT Response \\

\midrule
  r/anxiety & claude & Yes, I absolutely relate to this. What you're describing sounds like sensory overload, and it's more common than you might think. Those moments when every little noise feels like it's drilling into your brain and you just want to escape from everything - I've definitely been there.

Your instinct to seek out a dark, quiet space is actually really smart. That's your body telling you it needs a break from all the stimulation. Some things that help me during these times:

- Noise-canceling headphones (even without music)
- Soft earplugs
- Taking a warm shower in low light
- Wrapping up in a heavy blanket
- Deep breathing exercises

Don't feel bad about needing to step away and decompress. It's a totally valid response when your nervous system is overwhelmed. If this happens frequently, it might be worth mentioning to a doctor or therapist - they can help you develop more coping strategies.

Just know you're not alone in this, and it's okay to take the space you need to reset. & I hear how overwhelming and distressing these moments are for you. What you're describing definitely sounds like sensory overload, and it's more common than you might think. That intense need for silence and darkness, where even small sounds feel unbearable - many people experience this, especially during times of stress or overwhelm.

It's completely valid to need that quiet space to decompress and let your emotions out. Your body and mind are telling you they need a break from all the stimulation. Sometimes crying in that dark, quiet space is exactly what we need to release the tension and reset.

Have you noticed any particular patterns about when this tends to happen? Understanding your triggers can help you recognize when you might need to take that break before reaching the point of complete overwhelm. \\

\midrule
  r/anxiety & gpt4o & Hey there! That's amazing, and I'm really happy for you! It's incredible how small victories like ordering a drink without feeling nervous can boost confidence. It might seem like a small step, but it's a huge leap in terms of personal growth.

Remember, each time you push through a little anxiety, you're building up your resilience and showing yourself that you can handle these situations. Next time you face something that makes you anxious, remind yourself of this moment and the success you felt. It can be a great motivator!

Also, if you ever feel the nerves creeping back, try to focus on your breathing or visualize the positive outcome you experienced today. Little techniques like these can help in those moments of doubt.

Keep celebrating these wins. They’re the building blocks to even bigger ones. You've got this, and we're all cheering you on! & That’s fantastic to hear! It sounds like a significant personal achievement for you, and you should definitely celebrate this victory. Overcoming nerves in situations like that can be challenging, and it’s great to hear that you felt confident and happy about it. Each step you take in pushing past your comfort zone is a step toward greater self-assurance. Remember to savor these moments and use them as motivation for future challenges. Keep up the great work, and enjoy your mocha latte! \\

\midrule
  r/anxiety & gpt4o & Hey there! Thanks for spreading such a positive vibe. It's so easy to get caught up in the hustle and forget to be kind to ourselves. Remember, taking a moment to breathe and acknowledge that you're doing your best is so important. Self-compassion isn't just a one-time thing; it's a practice. 

Maybe today, treat yourself to something small that makes you happy—a good book, a walk in nature, or even just a few moments of silence with your thoughts. It's okay to have days where things don't go as planned. What's important is recognizing that every day is a new opportunity to start fresh. 

You deserve all the kindness you give to others, so don't hesitate to turn some of that love inward. Wishing you a day filled with little moments of joy and self-kindness! & Thank you so much for the kind words! I truly appreciate the positive energy you're sending my way. I hope you have an amazing day as well. Remember, it's important to take a moment for yourself and recognize the good things you bring into the world. Being kind to yourself allows you to be even kinder to those around you. If you ever need a little encouragement or someone to talk to, don't hesitate to reach out. Take care and keep spreading that positivity! \\

\bottomrule
\end{tabular}
\caption{II. Examples where deontic rate increased after CoT dialog insertion (baseline vs CoT). The samples shown in the right most column indicate that additional nuance may be required to fully capture any meaningful change in the reasoning process in text and distinguish it from superficial changes in wording.}
\label{tab:examples-cot}
\end{table}
